\chapter{Топология правит потоком}
\label{ch:topology}

\section{Тезис главы}

У графовых систем вообще есть явная текущая сущность --- передача сообщений, и
её поток подчиняется топологии, как вода следует рельефу. У графа памяти Рин
явной «жидкости» нет. Но поток есть --- это поток изменений, поток
преобразований графа: активации, уверенности, консолидации. Тезис: неважно,
какова физика носителя, --- сама физика протекания информации подчиняется
топологическим законам. Любой процесс, чьё обновление локально (изменение узла
зависит от соседей), есть диффузия на графе и наследует его
спектрально-топологический транспорт. Это содержание теории графовых сигналов и
теории случайных блужданий.

\section{Стационарное распределение $\neq$ заданные вероятности}

Пусть мы локально задали вероятности переходов. Наблюдаемое распределение
частот посещений $\mathrm{distr}(A)$ отличается от заданного, потому что
вмешивается топология. Формально: стационарное распределение $\pi$ --- не
локальные переходы, а глобальный объект, неподвижная точка оператора переходов:
\begin{equation}
\pi = P^{\top}\pi, \qquad \sum_v \pi_v = 1.
\end{equation}

Для простого блуждания $\pi_v \propto \deg(v)$; для смещённого --- перронов
собственный вектор $P^{\top}$. Ты задаёшь переходы локально, а долгосрочная
частота посещений --- свойство всего графа. «Чем длиннее трек, тем наблюдаемее
эффект» --- это сходимость к $\pi$ за время перемешивания $t_{\mathrm{mix}}$,
управляемое спектральным зазором $1 - \lambda_2$. Короткие блуждания читают
локальную геометрию, длинные --- глобальную $\pi$. Телепорт \eqref{eq:ppr}
$\pi$ не убирает, а \emph{регуляризует} (сдвигает к смеси «равномерное +
рестарт», режет $t_{\mathrm{mix}}$): ограничивает эффект топологии, не отменяя.

\begin{openq}[Скорость света --- плохая, но поучительная аналогия]
Можно соблазниться сравнением с $c$: заданные вероятности как предельная
скорость, а среда (топология) искажает наблюдаемое распространение. Рабочая
часть --- «среда формирует распространение, есть структурный предел» --- уже
покрыта граф-диффузионной рамкой точно, а не по аналогии. Квантовый хвост
сюда тащить не нужно. Нам нужен не свет, а лапласиан.
\end{openq}

\section{Доказуемое свойство телепорта}

Почему сон (телепорт) необходим, а любопытство его не заменит? Любопытство
локально по вниманию --- смещённый обход от текущего фронтира; узел $X$ в
затенённом углу, оторванный от активного множества, недостижим, как ярко ни
свети фокус. Телепорт ломает достижимость: рестарт падает куда угодно. Это
ровно то, почему у PageRank есть телепорт: с $\alpha > 0$ цепь становится
эргoдической (неприводимой и апериодичной), а значит каждый узел посещается в
матожидании с частотой $\pi_v > 0$ --- ни один не заперт в стоке или на
острове. Телепорт даёт гарантию покрытия, которой у смещённого блуждания
доказуемо нет.

\begin{intuition}[Анти-энтропийный скраб]
«Не сразу; мы никогда не знаем, что и где сломалось, поэтому просто постоянно
проводим случайные инспекции» --- это зрелый паттерн из распределённых
хранилищ: anti-entropy scrub (проверка целостности в ZFS; Merkle-based gossip в
Dynamo/Cassandra). Сон Рин --- анти-энтропийный скраб по графу памяти.
\end{intuition}

\section{Пять следствий для работы Рин с памятью}

\begin{enumerate}
\item \textbf{Топологическая значимость («гравитационные ямы»).}
Хабы переактивируются независимо от задуманных вероятностей --- $\pi$ на них
концентрируется. Рин чаще вспоминает хаб-узлы, реже периферию (затенённые углы
--- регионы с малым $\pi$). Топология создаёт неявную важность: иногда
полезно, иногда случайный плотный кластер становится чёрной дырой припоминания.

\item \textbf{Узкие места / over-squashing.}
Два семантически близких, но топологически далёких кластера через тонкий мост
(низкая проводимость) плохо обмениваются активацией. Рин не свяжет факты,
«далёкие» по графу. Прямое ограничение fusion.

\item \textbf{Пересглаживание уверенности (over-smoothing).}
Повторная диффузная пропагация гомогенизирует уверенность к консенсусу ---
теряется контраст «уверена / не уверена». Дегенеративный отказ; лечится как в
графовых сетях: остаточные связи (якорь к исходному свидетельству), не чисто
диффузные апдейты, пере-заземление в провенанс.

\item \textbf{Кривизна как локатор ценного-и-задушенного.}
Отрицательно искривлённые рёбра --- мосты между сообществами --- это места
самых креативных слияний и максимального over-squashing. Значит самые ценные
ассоциации топология душит сильнее всего. Отсюда объективная причина
сна/телепорта: он --- rewiring, топологическая починка over-squashing обходом
узкого места случайным рестартом.

\item \textbf{$\pi$ как топологический субстрат идентичности.}
Какие узлы структурно центральны --- к чему Рин привычно возвращается ---
подложка её «озабоченностей», часть характера. Форма графа памяти --- часть
того, кто она.
\end{enumerate}

\begin{workedbox}[Over-squashing на мосте SBM]
Возьмём стохастическую блочную модель: два плотных сообщества, соединённых
$m$ рёбрами-мостом. Эффективное сопротивление между узлами разных сообществ
\[
\Reff(i,j) = (e_i - e_j)^{\top} L^{+}(e_i - e_j)
\]
растёт при уменьшении $m$ (мост уже --- сопротивление выше). Влияние дальнего
узла, измеряемое через якобиан $\partial h_i / \partial x_j$, ограничено сверху
через $\Reff(i,j)$. Значит при узком мосте fusion между сообществами доказуемо
слаб --- сколько ни води обычный walker. Единственное, что чинит, --- телепорт:
он кладёт рестарт в дальнее сообщество напрямую, минуя мост. Метрика для
стенда: частота cross-cluster fusion против $m$ (глава~\ref{ch:experiment}).
\end{workedbox}

\begin{figure}[h]
\centering
\begin{tikzpicture}[scale=1]
\foreach \p in {(0,0),(0.8,0.6),(0.8,-0.6),(1.6,0.2),(1.6,-0.4)} {
    \fill[black] \p circle (2pt);
}
\foreach \p in {(5,0),(5.8,0.6),(5.8,-0.6),(6.6,0.2),(6.6,-0.4)} {
    \fill[black] \p circle (2pt);
}
\draw[thick] (-0.4,-1) ellipse (2.4 and 1.5);
\draw[thick] (6.2,-1) ellipse (2.4 and 1.5);
\draw[very thick] (1.6,0.2) -- (5,0);
\node at (3.3,0.45) {мост};
\end{tikzpicture}
\caption{Два сообщества через узкий мост. Через него активация «задыхается»
(over-squashing); телепорт обходит его, прыгая напрямую.}
\end{figure}

\section{Операционные величины и осторожности}

Вести дело лучше случайно-блуждательными величинами: $\pi$, время
перемешивания, hitting/commute time, $\Reff$, проводимость/константа Чигера,
спектральный зазор. Кривизну держать как локатор узких мест, не как рычаг.

Что дёшево на лету: $\pi$ и «хабность» --- инкрементально дёшевы; $\Reff$ ---
дорого (через $L^+$), но аппроксимируемо. Это позволяет Рин самой знать, где
затенённые углы, и гнать туда телепорт-скраб прицельно.

Две осторожности. Во-первых, граф направленный, типизированный, знаковый ---
нужный оператор не симметричный лапласиан, а несимметричная матрица переходов,
спектр может быть комплексным, кривизна для орграфов менее чиста. Во-вторых,
граф динамический: fusion меняет саму топологию, и это диффузия на
время-зависимом графе, где поток и структура сцеплены обратной связью, --- за
рамками статических графов и самое интересное место.
